\section*{2. 比较定理}

\subsubsection*{2.1.}
对每个 $\mathbb C$ 上的有限型概形 $X$，用 $X_{\mathrm{cl}}$ 表示赋予通常拓扑的空间 $X(\mathbb C)$。

设 $S$ 是 $\mathbb C$ 上的光滑曲线，即 $\mathbb C$ 上纯 $1$ 维的光滑分离有限型概形。设 $0\in S(\mathbb C)$，$f:X\longrightarrow S$ 是分离有限型 $S$-概形，$\Lambda$ 是 Noether 挠环（例如 $\mathbb Z/n\mathbb Z$）。令 $X_0=f^{-1}(0)$。

\subsubsection*{2.2.}
用 $S^h$ 表示 $S$ 在 $0$ 处 Hensel 化所得的 Hensel 性离散赋值谱，用 $\eta$ 表示 $S^h$ 的一般点。有一个函子
\[
\begin{aligned}
\Psi:(\text{层，底空间为 }X)
&\xrightarrow{\ \mathrm{restriction}\ }
(\text{层，底空间为 }X\times_SS^h)\\
&\xrightarrow{\ \Psi\ }
(\text{层，底空间为 }X_0\times S^h).
\end{aligned}
\]
其导出函子为复合
\[
R\Psi:D(X,\Lambda)
\longrightarrow D(X\times_SS^h,\Lambda)
\xrightarrow{\ R\Psi\ }D(X_0\times S^h,\Lambda).
\]

\subsubsection*{2.3.}
设 $D\subset S_{\mathrm{cl}}$ 是包含 $0$ 的圆盘。复合函子 $\Psi_{\mathrm{cl}}$（或简记 $\Psi$）
\[
\begin{aligned}
\Psi_{\mathrm{cl}}:(\text{层，底空间为 }X_{\mathrm{cl}})
&\xrightarrow{\ \mathrm{restriction}\ }
(\text{层，底空间为 }X_{\mathrm{cl}}\times_{S_{\mathrm{cl}}}D)\\
&\xrightarrow{\ \Psi\ }
(\text{层，底空间为 }(X_0)_{\mathrm{cl}}\times\mathcal D)
\end{aligned}
\]
在唯一同构意义下不依赖于 $D$ 的选择；回忆一下，将 $D$ 换为较小圆盘时，$\mathcal D$ 不变。其导出函子为复合
\[
R\Psi_{\mathrm{cl}}:D(X_{\mathrm{cl}},\Lambda)
\longrightarrow D(f_{\mathrm{cl}}^{-1}(D),\Lambda)
\xrightarrow{\ R\Psi\ }
D((X_0)_{\mathrm{cl}}\times\mathcal D,\Lambda).
\]

\subsubsection*{2.4.}
同样定义
\[
\begin{aligned}
R\Psi_\eta &:D(X-X_0,\Lambda)
  \longrightarrow D(X_0\times\eta,\Lambda),\\
R\Psi_{\eta\,\mathrm{cl}} &:D((X-X_0)_{\mathrm{cl}},\Lambda)
  \longrightarrow D((X_0)_{\mathrm{cl}}\times\mathcal D^*,\Lambda).
\end{aligned}
\]

\subsubsection*{2.5.}
1.1.11 的构造定义了从 $\mathcal D^*$ 到 $\eta$ 的拓扑斯态射，它延拓为
\[
\varepsilon:\mathcal D\longrightarrow S^h.
\]
结合 $\varepsilon:(X_0)_{\mathrm{cl}}\longrightarrow X_0$（SGA 4, XI, 4），此态射定义
\[
\varepsilon:(X_0)_{\mathrm{cl}}\times\mathcal D
\longrightarrow X_0\times S^h.
\]
若 $\widetilde D^*$ 是 $D^*$ 的一个万有覆盖，而 $k(\bar\eta)$ 是 $k(\eta)$ 的相应代数闭包（1.1.10），则
\[
\varepsilon^*(F_0,F_\eta,\alpha)
=\bigl(\varepsilon^*F_0,\varepsilon^*F_{\bar\eta},
\varepsilon^*\alpha\bigr).
\]

